\section{Permanence from the cut-off down to a fixed number of rows}\label{meso}
The goal of this section is Proposition~\ref{prop:salida-banda-meso}: under $\nu$, a history with
regular entry at $M$ has its first exit at a size in the band $[2^jK,2^{j+1}K)$ with probability at
most $Ce^{-c\,2^jK}$. To this end we control, between $M$ and $K$, the number of empty crossings on
each diagonal (Lemmas~\ref{meso:prescripciones} and~\ref{lema:presupuesto-meso}) and certain
observables of rows, columns, fibers, and rectangles (Proposition~\ref{prop:concentracion-meso}).
From them we reconstruct the Bethe point and check that the board remains in the inner class
(Proposition~\ref{meso:reconstruccion}). The order in which the constants are fixed is given in
Section~\ref{meso:constantes}.

We use $\mathcal F^{(k)}=\mathcal F_{n-k}$ for the history through the board $B^{(k)}$ of size $k$.
Let $M=\lceil\delta n\rceil$. For
the board of size $k$, write $D_{\max}(k)$ for the maximum, over diagonals of both families of the initial
board, of the number of empty crossings on the line, whether available or not. This quantity is
nonincreasing along the history and bounds from above the number of available squares on any diagonal in
either family. All probabilities in this section are under the adaptive law $\nu$ started at the initial
board. Let $G_k^-$ denote the event that, upon reaching size $k$, all boards of strictly larger size have
been regular (Definition~\ref{def:regular}). We include in $G_k^-$ the transition to the first-exit board.

\subsection{Prescriptions and the diagonal budget}
\begin{lema}[One-step bounds]\label{meso:paso}
At a regular parent of size $k$, let $\nu_B(j\mid i)$ be the adaptive probability of choosing column $j$
given that row $i$ is chosen, with $\nu_B(j\mid i)=0$ on unavailable squares. There are constants
$C_q,C_{\rm c}$, which can be fixed before $\delta$ and the diagonal mass bound $\varepsilon$, such that,
for all sufficiently large $n$ and all $k$ above a fixed threshold,
\begin{align*}
 \nu(i\text{ chosen}\mid\mathcal F^{(k)})&=1/k,\\
 \left|\nu(j\text{ chosen}\mid\mathcal F^{(k)})-1/k\right|&\le C_{\rm c}/k^2,\\
 \nu((i,j)\text{ chosen}\mid\mathcal F^{(k)})&\le C_q/k^2 .
\end{align*}
\end{lema}
\begin{proof}
For a regular parent, the child expansion of Theorem~7.1 of~\cite{pred} and of the proof of its
Theorem~8.1 \cite[Sections~7.8, 8.2 and~8.5]{pred}, which is item (I2) of Section~\ref{leyes}, gives
$e^{L(B+e)-L(B)}=p_e(1+d_e+O(k^{-2}))$, with $d_e=O(k^{-1})$, uniformly over its available squares.
This includes a child at the first exit. With $B_{\rm RC}$ the row--column incidence matrix, the exact
constraint $B_{\rm RC}(p\odot d)=0$ of \cite[Section~6.1]{pred}, which is item (I3), gives both
$\sum_jp_{ij}d_{ij}=0$ and $\sum_i p_{ij}d_{ij}=0$. Since each row has unit $p$-mass, its normalizing
factor is $1+O(k^{-2})$. Hence on every available square the row-normalized law satisfies
\[
 \nu_B(j\mid i)=p_{ij}\bigl(1+d_{ij}+O(k^{-2})\bigr).
\]
Summing in $i$, the column identity cancels the linear term. The errors sum to $O(k^{-2})$ because
$\sum_i p_{ij}=1$. Thus $\sum_i\nu_B(j\mid i)=1+O(k^{-2})$, and the uniform choice of row gives
$\nu(j\text{ chosen}\mid\mathcal F^{(k)})=1/k+O(k^{-3})$. In what follows the weaker bound of the
statement suffices.

For $k>M$ the last remainder is $O_\delta(k^{-3})$. Once $\delta$ is fixed, increasing $n$ bounds
it by $k^{-2}$ with coefficient one and gives $\nu_B(j\mid i)\le2p_{ij}$. The diffusion constant for
these parents comes from the profile. For $k\le M$ the constants in the expansion depend only on the
fixed diffusion, vacancy, and gap bounds of the outer class, provided the diagonal mass is below the
threshold of \cite[Theorem~8.1]{pred}. These two estimates fix $C_q,C_{\rm c}$ before $\delta$ and
$\varepsilon$.
\end{proof}

\begin{lema}[Prescribed rows and columns]\label{meso:prescripciones}
For any fixed $s$ rows and $t$ columns of the initial board and every size $m$ with
$K_0\le m\le n$,
\begin{equation}\label{prescr}
 \nu(\text{the $s$ rows and $t$ columns remain empty at size $m$},G_m^-)
 \le (m/n)^{s+t}\exp\{C(st+s+t)/m\}.
\end{equation}
If the prescriptions cannot fit, the left-hand side is zero. The same constant $C$ applies to all values of
$s,t$.
\end{lema}
\begin{proof}
At a regular parent of size $k$ containing the prescribed indices, Lemma~\ref{meso:paso} gives probability
$s/k$ of choosing a prescribed row, probability at least $t/k-C_{\rm c}t/k^2$ of choosing a prescribed
column, and probability at most $C_qst/k^2$ for their intersection. Within each family the events are
disjoint. The survival factor is thus at most $\exp[-(s+t)/k+C_1(st+s+t)/k^2]$. We iterate from $n$ down to
$m+1$, discarding histories that already had a non-regular parent. The inequalities $\sum_{k=m+1}^n
k^{-1}\ge\log(n/m)-1/m$ and $\sum_{k=m+1}^n k^{-2}\le1/m$ give \eqref{prescr}, even when $s,t$ depend on
$m$.
\end{proof}

Fixing $s$ squares on a diagonal prescribes exactly $s$ distinct rows and $s$ distinct columns, regardless
of availability. There are fewer than $4n$ diagonals, each allowing at most $\binom ns$ choices. By
\eqref{prescr} and $\binom ns\le(en/s)^s$, there is a fixed $C_{\rm geo}<\infty$ such that, for $1\le s\le
m/4$,
\begin{equation}\label{geompres}
 \nu(D_{\max}(m)\ge s,G_m^-)\le 4n\left(\frac{C_{\rm geo}m^2}{ns}\right)^s .
\end{equation}
Indeed, the factor $\exp[C(s^2+2s)/m]$ can be absorbed into a constant raised to the power $s$. For $s>m$
the event is empty.
